%% kapa: exam=1 round=36 year=2025 period=2 form=A subject=accounting subject_ko=회계학 seq=5 number=45
%% source: Q-Net 감정평가사 기출문제 (cst00309) artlSeq=5231525 file=raw/1008/5231525/2025년 제36회 감정평가사 1차시험 2교시 문제지 원본.pdf
%% conversion: kordoc 4.24.51 -> tools/gen1.py (auto); review=proofread-sonnet; flags=none
\begin{question}{45}
(주)감평은 20×1년 5월 초 공장건물 신축공사를 시작하여 20×2년 10월 말 완공하였다. 동 건물은 차입원가 자본화 대상인 적격자산이다. 동 건물 신축 관련 자료가 다음과 같을 때, (주)감평이 20×1년에 자본화할 차입원가는? (단, 연평균지출액과 이자비용은 월할 계산한다.)
\begin{tcolorbox}[enhanced, colback=white, colframe=black, boxrule=0.4pt, arc=0pt, left=4pt, right=4pt, top=4pt, bottom=4pt]
(1) 공장건물 신축을 위한 공사비 지출내역은 다음과 같다. 지출내역 중 20×1년 9월 1일에 지출한 ￦108,000은 전액 공장건물 신축과 관련하여 동 일자에 수령한 정부보조금(상환의무 없음)으로 지출되었다.
\begin{center}\small
\begin{tabular}{c@{\hspace{3em}}r}
\uline{\ 일\ 자\ } & \uline{\ 금\ 액\ } \\
20×1. 5. 1. & ￦240,000 \\
20×1. 9. 1. & 108,000 \\
20×2. 6. 1. & 180,000
\end{tabular}
\end{center}
(2) 일반차입금 현황은 다음과 같으며, 모든 차입금의 이자는 단리로 계산하여 1년 후급하는 조건이다. 또한 특정차입금은 없다.
\begin{center}\small
\begin{tabular}{c r c c}
\uline{\ 구\ 분\ } & \uline{\ 차입금액\ } & \uline{\hspace{3em}차입기간\hspace{3em}} & \uline{\ 연 이자율\ } \\
차입금 A & ￦100,000 & 20×0. 1. 1.∼20×2. 12. 31. & 12 \% \\
차입금 B & 200,000 & 20×1. 7. 1.∼20×2. 06. 30. & 10 \%
\end{tabular}
\end{center}
\end{tcolorbox}
\begin{choices}
  \item ￦17,600
  \item ￦18,860
  \item ￦19,440
  \item ￦21,560
  \item ￦22,000
\end{choices}
\end{question}
\answer{1}
